\documentclass[10pt,a4paper]{amsart}
\usepackage[margin=19mm]{geometry}
\usepackage{amsmath,amssymb,mathtools}
\usepackage[scr=rsfs,cal=euler]{mathalfa}
\usepackage{xfrac}
\usepackage[unicode,hidelinks,bookmarks=false]{hyperref}
\newcommand{\ie}{i.e.\ }
\newcommand{\cf}{cf.\ }
\newcommand{\resp}{respectively\ }
\newcommand{\Subsection}{\subsection}
\providecommand{\nobreakdash}{\nobreak\hbox{-}}
\setlength{\emergencystretch}{2em}
\multlinegap0pt
\usepackage{bm}
\usepackage{bbm}
\usepackage{dsfont}
\usepackage{xspace}
\usepackage{cancel}
\usepackage{mathtools}
\usepackage[shortlabels]{enumitem}
\usepackage{amsthm}
\makeatletter
\@ifundefined{definition}{\newtheorem{definition}{Definition}}{}
\@ifundefined{lemma}{\newtheorem{lemma}{Lemma}}{}
\@ifundefined{theorem}{\newtheorem{theorem}{Theorem}}{}
\makeatother
\def\r{\mathbb{R}}
\title[Exponential dichotomy and $(L^p,L^q)$-admissibility]{Exponential dichotomy and $(L^p,L^q)$-admissibility}
\author{Trinh Viet Duoc, Nguyen Van Trong}
\date{Source: arXiv:2601.06534v1 (CC BY 4.0)}
\begin{document}
\maketitle

\noindent\emph{Source and licence.} Exponential dichotomy and $(L^p,L^q)$-admissibility. Version arXiv:2601.06534v1 (CC BY 4.0), distributed under the Creative Commons Attribution 4.0 International licence, as recorded in the arXiv licence metadata for this exact version and shown on the abstract page https://arxiv.org/abs/2601.06534v1. Full rights notice: licenses/arxiv-2601.06534-CC-BY-4.0.txt.\par\smallskip

\noindent\textit{Editorial scope.} Counted unit: the Theorem whose source label is \texttt{thm:2.2} in the article (source0.tex lines 267--269, source label \texttt{thm:2.2}), delivered together with the article's own complete original proof of it (source0.tex lines 270--339, one \texttt{proof} environment of 70 lines). No mathematics is written, repaired or re-derived: statement and proof are the article's own text line for line. Required context retained uncounted: 2.6 (lines 139--158); 2.7 (lines 189--191); lem1 (lines 193--195). Every definition, display and label of those lines is retained unchanged. Omitted from this package and not redistributed: 1--138, 159--188, 192--192, 196--266, 340--664. Nothing inside a retained excerpt is omitted or altered: every retained body is delivered exactly as the article's lines are written, so no omission receipt of any kind accompanies this package. Citations: the retained excerpts carry no citation command at all, so no bibliography entry of the article is transcribed and no reference is redistributed. Reference adaptation: none was needed and none was made; every reference inside the delivered unit resolves inside it, so no unresolved reference or citation is produced. Printed slips kept literal: no printed word, symbol, hypothesis or number of the counted statement or of its proof was altered. Layout and class: the article's own class is \texttt{amsart} with packages including latexsym, amsmath, amsfonts, amscd, amssymb, verbatim, mathtools, enumerate, xcolor, amsthm, csquotes; the wrapper is the lane's pinned \texttt{amsart} editorial wrapper at 10pt on A4, which restates the theorem-like environments the retained text needs and adds the article's own macro definitions that the retained text needs (\texttt{\textbackslash r}), with extra wrapper packages (bm, bbm, dsfont, xspace, cancel, mathtools, enumitem). The delivered numbering is the unit's own. Assets: no retained span contains a figure environment, a graphics-inclusion command, a PDF-page inclusion, a table or a TikZ picture; no image file is copied into this package, the package declares no asset dependency and no image file is redistributed with it. Licence: version arXiv:2601.06534v1 of this article is distributed under the Creative Commons Attribution 4.0 International licence, as recorded in the arXiv licence metadata for this exact version and shown on the abstract page https://arxiv.org/abs/2601.06534v1. A full rights notice is delivered at licenses/arxiv-2601.06534-CC-BY-4.0.txt. Characters: every retained excerpt is pure ASCII, so the delivered engine, XeLaTeX with its default Latin Modern text font, reports no missing character.
\begin{definition}
An evolutionary family $\left( T(t,\tau)\right)_{t \geq \tau}$ with respect to a family of norms $\Vert \cdot\Vert _t$ on the line is said to be an exponential dichotomy with respect to the family of norms $\Vert \cdot\Vert _t$ if it satisfies the following conditions:
\begin{enumerate}[i.]
\item There exist dichotomic projections $P(t)$ in $B(X), t \in \mathbb{R}$ such that 
\begin{equation*}
P(t)T(t,\tau)=T(t,\tau)P(\tau) \quad \text{ for } t\geq \tau. 
\end{equation*}
\item The restriction mapping  
$T(t,\tau): Ker P(\tau) \rightarrow Ker P(t) $
is an isomorphism for $t \geq \tau$. We denote its inverse $(T(t,\tau)_{|_{\text{\rm Ker}P(\tau)}})^{-1}$ by $T(\tau,t)_|$.
\item There exist constants $a<0<b$ and $D>0$ such that
\begin{align}
\begin{split}
\Vert T(t,\tau)P(\tau)x \Vert _{\tau} &\leq De^{a(t-\tau)}\Vert x\Vert _{\tau},\\
\Vert T(\tau,t)_| Q(t)x \Vert _{t} &\leq De^{-b(t-\tau)}\Vert x\Vert _{t},
\end{split} \label{2.6}
\end{align}
for $t \geq \tau, x\in X$ and $Q(t)=I-P(t)$.
\end{enumerate}
\end{definition}

\begin{equation}
x(t)=T(t,\tau)x(\tau)+\displaystyle\int\limits_{\tau}^{t}T(t,s)y(s)ds  \quad  \text{for} \quad t\geq \tau. \label{2.7}
\end{equation}

\begin{lemma}\label{lem1}
Given $y\in Y_2$. If a function $x: \r \to X$ satisfies \eqref{2.7}, then $x$ is continuous on $\r$.
\end{lemma}

\begin{theorem}\label{thm:2.2}
If the evolutionary family $(T(t,\tau))_{t \geq \tau}$ is exponentially dichotomic with respect to a family of norms $\Vert \cdot\Vert _t$, then for each $y \in Y_2$ there exists a unique $x \in Y_1$ satisfying \eqref{2.7}.
\end{theorem}

\begin{proof}
For $t\in \mathbb{R}$, we define
$$ x_1(t)=\displaystyle\int\limits_{-\infty}^{t}T(t,\tau)P(\tau)y(\tau)d\tau \quad \text{and}\quad  x_2(t)=\displaystyle\int\limits_{t}^{\infty}T(t,\tau)_|Q(\tau)y(\tau)d\tau.$$
Using Holder's inequality and \eqref{2.6}, we have
\begin{align*}
\|x_1(t)\| &\le \Vert x_1(t)\Vert _t  \le \displaystyle \int_{-\infty}^{t} \Vert T(t,\tau)P(\tau)y(\tau)\Vert _t d\tau  
\le \displaystyle \int_{-\infty}^{t} D e^{{a}(t-\tau )}\Vert y(\tau)\Vert _{\tau} d\tau \nonumber \\
&= \displaystyle\sum_{m=0}^{\infty} \displaystyle \int_{t-m-1}^{t-m} D e^{{a}(t-\tau )}\Vert y(\tau)\Vert _{\tau} d\tau 
\leq \displaystyle\sum_{m=0}^{\infty}D e^{{a}m} \displaystyle \int_{t-m-1}^{t-m} \Vert y(\tau)\Vert _{\tau} d\tau \nonumber  \\
&\le \displaystyle\sum_{m=0}^{\infty}D e^{{a}m} \left( \displaystyle \int_{t-m-1}^{t-m} \Vert y(\tau)\Vert _{\tau}^q d\tau \right)^{1/q} 
\leq \displaystyle\sum_{m=0}^{\infty}D e^{{a}m} \Vert y\Vert_q 
=\frac{D}{1-e^{a}}\Vert y\Vert _q, \\
\|x_2(t)\| &\le \Vert x_2(t)\Vert _t  \le \displaystyle \int_{t}^{+\infty} \Vert T(t,\tau)_|Q(\tau)y(\tau)\Vert _t d\tau  
\le \displaystyle \int_{t}^{+\infty} D e^{-{b}(\tau -t)}\Vert y(\tau)\Vert _{\tau} d\tau \nonumber \\
&= \displaystyle\sum_{m=0}^{\infty} \displaystyle \int_{t+m}^{t+m+1} D e^{-{b}(\tau -t)}\Vert y(\tau)\Vert _{\tau} d\tau 
\leq \displaystyle\sum_{m=0}^{\infty}D e^{-{b}m} \displaystyle \int_{t+m}^{t+m+1} \Vert y(\tau)\Vert _{\tau} d\tau \nonumber  \\
&\le \displaystyle\sum_{m=0}^{\infty}D e^{-{b}m} \left( \displaystyle \int_{t+m}^{t+m+1} \Vert y(\tau)\Vert _{\tau}^q d\tau \right)^{1/q} 
\leq \displaystyle\sum_{m=0}^{\infty}D e^{-{b}m} \Vert y\Vert_q 
=\frac{D}{1-e^{-{b}}}\Vert y\Vert _q. \notag
\end{align*}
Therefore, $x_1$ and $x_2$ are well-defined. Furthermore,
\begin{equation}\label{eq1}
 \|x_1(t)-x_2(t)\|_t \le \Big(\frac{D}{1-e^{a}} + \frac{D}{1-e^{-{b}}}\Big) \Vert y\Vert _q \quad \text{for all }\, t\in \r.
\end{equation}
Acting $P(t)$ to the both sides of \eqref{2.7}, we have
$$ P(t)x(t)= T(t,\tau)P(\tau)x(\tau)+\displaystyle\int\limits_{\tau}^{t}T(t,s)y(s)ds  \quad  \text{for} \quad t\geq \tau.$$
For $\tau \to -\infty$, we obtain
$$ \lim_{\tau \to -\infty}T(t,\tau)P(\tau)x(\tau) =  P(t)x(t) - x_1(t).$$
We have
$$\|T(t,\tau)P(\tau)x(\tau)\| \le \|T(t,\tau)P(\tau)x(\tau)\|_{t} \le D e^{a(t-\tau)} \|x(\tau)\|_{\tau} \le D e^{a(t-\tau)} \|x\|_{\infty}.$$
This leads to $P(t)x(t) = x_1(t)$ for each $t \in \mathbb{R}$.

Fix $\tau \in \mathbb{R} $, applying $Q(t)$ to the both sides of \eqref{2.7}, we have
$$ Q(t)x(t)= T(t,\tau)Q(\tau)x(\tau)+\displaystyle\int\limits_{\tau}^{t}T(t,s)Q(s)y(s)ds  \quad  \text{for} \quad t\geq \tau.$$
By $T(t,\tau): Q(\tau)X \to Q(t)X$ is an isomorphism, 
$$ T(\tau,t)_| Q(t)x(t) = Q(\tau)x(\tau)+\displaystyle\int\limits_{\tau}^{t}T(\tau,s)_|Q(s)y(s)ds \quad  \text{for} \quad t\geq \tau.$$
For $t \to +\infty$, we obtain
$$ \lim_{t\to +\infty}T(\tau,t)_| Q(t)x(t) = Q(\tau)x(\tau)+ x_2(\tau).$$
Same as above, we get $Q(\tau)x(\tau)= -x_2(\tau)$ for each $\tau \in \mathbb{R}$.

So, we have $x(t)=x_1(t)-x_2(t)$ for all $t\in \mathbb{R}$. Therefore, $x\in Y_1$ satisfying \eqref{2.7} is unique.
Next, we show that the functions  $x_1$ and $x_2$ belong to $L^p$.
We only consider $x_2$ since the argument for $x_1$ is analogous. 
Since $p \geq q$, there exists $r \in [1,\infty]$ such that
$$\dfrac{1}{r}+\dfrac{1}{q}=1+\dfrac{1}{p}.$$ 
Let
$$
 \omega (t)=
 \begin{cases}
 0& \text{if}\quad t \geq 0,\\
e^{{b}t}& \text{if}\quad t<0\\ 
 \end{cases}
 $$
 and $z(t)=\Vert y(t)\Vert_t$. We have $\omega \in L^{r}(\r)$ and $z\in L^q(\mathbb{R})$.
By Young's inequality  for convolution,   $w*z \in L^p(\mathbb{R})$. Furthermore,
$$(\omega *z)(t)= \displaystyle \int _{-\infty}^{+\infty}\omega (t-\tau)z(\tau)d\tau=\displaystyle \int _{t}^{+\infty}e^{-{b}(\tau -t)}\Vert y(\tau) \Vert_{\tau}d\tau.$$
We also have $\Vert x_2(t)\Vert _t \le D(\omega *z)(t)$ for all $t$. Thus, $x_2 \in L^p$ and
$$\Vert x_2\Vert_p \leq D \Vert \omega *z\Vert _p \le \frac{D}{(br)^{\frac{1}{r}}} \|y\|_q.$$
Similarly, $x_1 \in L^p$. Therefore, $x_1 - x_2 \in L^p$ and
$$ \Vert x_1 -x_2 \Vert_p \le \Big(\frac{D}{(-ar)^{\frac{1}{r}}}+ \frac{D}{(br)^{\frac{1}{r}}}\Big) \|y\|_q.$$
Let $x(t)=x_1(t)-x_2(t)$ for $t \in \mathbb{R}$, and take $t_0 \in \mathbb{R}$. We now show that $x$ satisfies \eqref{2.7} for $t\ge t_0$. 
\begin{align}
x(t)&=\displaystyle\int\limits_{t_0}^{t}T(t,\tau)y(\tau)d\tau-\displaystyle\int\limits_{t_0}^{t}T(t,\tau)P(\tau)y(\tau)d\tau -\displaystyle\int\limits_{t_0}^{t}T(t,\tau)Q(\tau)y(\tau)d\tau \nonumber
\\&+\displaystyle\int\limits_{-\infty}^{t}T(t,\tau)P(\tau)y(\tau)d\tau-\displaystyle\int\limits_{t}^{+\infty}T(t,\tau)_|Q(\tau)y(\tau)d\tau \nonumber\\
&=\displaystyle\int\limits_{t_0}^{t}T(t,\tau)y(\tau)d\tau+\displaystyle\int\limits_{-\infty}^{t_0}T(t,t_0)T(t_0,\tau)P(\tau)y(\tau)d\tau-\displaystyle\int\limits_{t_0}^{+\infty}T(t,t_0)T(t_0,\tau)_|Q(\tau)y(\tau)d\tau \nonumber\\
&=\displaystyle\int\limits_{t_0}^{t}T(t,\tau)y(\tau)d\tau+T(t,t_0)x(t_0).\nonumber 
\end{align}
By Lemma \ref{lem1} and \eqref{eq1}, $x\in C_b$.
So, there exists a unique $x \in Y_1$ satisfying \eqref{2.7}.
\end{proof}

\end{document}
